\section{Transport and Delivery}\label{sec:transport}

Above the radio the question is the same: deliver each frame before its deadline, and send less when the
path cannot.

\subsection{Earlier congestion signals}
Delay- and loss-based controllers react after a queue has formed, which is a frame too late for XR.
\emph{L4S} marks congestion before queues grow. In volumetric video conferencing, an L4S pipeline
converged on the available rate 45\% faster than Google Congestion Control (GCC) and lost no packets,
while the WebRTC baseline left 3\% of frames undecodable \cite{W4408870658}. \emph{Comparisons of
existing controllers} remain mixed: integrated into ALVR, GCC tracked the available rate closely with few
dropped frames, but only NADA held a fair share against loss-based flows \cite{W4415538381}.
\emph{In-network adaptation} moves the decision to the base station. Octopus lets senders overestimate
capacity and drops application-chosen packets at the cell, reporting 1.5--18$\times$ better performance
for standard and volumetric video \cite{W4401386494}.

\subsection{Transports with deadlines}
QUIC removes transport head-of-line blocking and offers unreliable datagrams. \emph{Media over QUIC} (MoQ)
adds delivery timeouts, which a point-cloud server used for implicit adaptation: clients with tight
latency targets received lower quality, relaxed ones higher quality \cite{W4417496773}. WebRTC keeps its
strengths for groups: through a selective forwarding unit, volumetric video adapted to user movement 13\%
and 36\% faster than over TCP at 5 and 10~ms of network latency \cite{W7151958829}. Which transport suits
interactive XR is still \emph{contested}.

\Cref{sec:qoe} showed that loss, not delay, drives QoE, yet a retransmission costs a frame time.
\emph{Selective protection} spends redundancy where it matters. ReVo applies forward error correction (FEC)
to critical colour and depth content and repairs the rest with a neural module, cutting freezes by up to
95.7\% on real loss traces \cite{W7159678172}. A hybrid broadcast--unicast design masked typical wireless
loss with 15\% FEC overhead while keeping volumetric latency under 40~ms \cite{W4415933998}.

\subsection{Adapting what is rendered and sent}
\emph{Application-layer rate adaptation} is the oldest loop and still effective: over Wi-Fi, a
network-aware bitrate algorithm for ALVR adapted to contention and mobility \cite{W4405014964}.
\emph{Perception-driven sending} lets the user decide what to send, and it is the one family validated
with users. Dynamic foveation driven by streamed gaze added 0.60 MOS over static foveation while saving
9.81\% bitrate in cloud VR gaming \cite{W4377981508}, and EyeNexus adapted the foveal region to bandwidth
and cut latency by up to 70.9\% \cite{W4416659623}. Encoding settings driven by a learned QoE model raised
interaction quality by 1.20 MOS \cite{W4400981668}. \emph{Split decisions} move work between device and
edge: RenderFusion switches objects between local and remote rendering using a perceptual study of
latency \cite{W4389296467}.

\subsection{Volumetric and Gaussian-splat delivery}
Volumetric content turns the trade-off between bandwidth and latency into one between compression and
latency. In VR2Gather, compression cut bandwidth by more than 90\% but raised latency to 278.7~ms, and
DASH delivery reached 562~ms against 87.4~ms over TCP \cite{W4362496249}. MuV2 served over 30 untethered devices from
one access point by mixing direct volumetric streaming with remote rendering \cite{W4399121637}. Gaussian
splats bring compact, progressive formats: V\textsuperscript{3} streams about 0.5~MB per frame and renders
at 27~fps on a phone \cite{W4404525293}, Vega reaches 30~fps on smartphones by prioritising important
objects \cite{W4416477288}, and TIGAS renders splats remotely and streams 2D views over QUIC
\cite{W7160906632}. These systems report PSNR, SSIM and frame rate; none of them reports a networked user
study.
